\section{Introduction}
The famous \emph{Plünnecke--Ruzsa} inequalities state that if $K\ge1$ is a parameter and $A$ is a finite subset of an abelian group satisfying the \emph{small doubling} hypothesis $|A+A|\le K|A|$, then $|mA-nA|\le K^{m+n}|A|$ for all non-negative integers $m,n$ \cite{plunnecke,ruz.pl.1,ruz.pl.2} (see also \cite{petridis} for a dramatically simplified proof). It has long been known that the directly analogous statement fails for arbitrary non-abelian groups, and that a bound of the form $|A^2|\le K|A|$ does not in general imply a bound of the form $|A^3|\le O_K(|A|)$ (see e.g. \cite[Example 2.5.2]{book}, which is essentially reproduced from \cite{tao.product.set}). Nonetheless, under the stronger \emph{small tripling} assumption $|A^3|\le K|A|$, we do have that $|A^{\epsilon_1}\cdots A^{\epsilon_m}|\le K^{3(m-2)}|A|$ for all choices of $\epsilon_i=\pm1$ and all $m\ge3$ \cite[Proposition 2.5.3]{book}. A similar result holds more generally if $A$ is an open precompact subset of a locally compact group, with Haar measure in place of cardinality \cite[Lemma 3.4]{tao.product.set}.

Breuillard and the second author showed that if $A$ is a \emph{ball} in a Cayley graph then the small-tripling hypothesis can \emph{almost} be weakened to a small-doubling hypothesis, as follows.
\begin{theorem}[Breuillard--Tointon {\cite[Lemma 2.10]{bt}}]\label{lem:bt}
Let $K\ge1$. Suppose $S$ is a finite symmetric subset of a group containing the identity, and $|S^{2n+1}|\le K|S^n|$ for some $n\in\N$. Then $|S^{3n}|\le O_K(|S^n|)$, and $S^{2n}$ is an $O_K(1)$-approximate group.
\end{theorem}
A \emph{$K$-approximate group} is a symmetric subset $A$ of a group containing the identity such that $A^2$ is covered by at most $K$ left translates of $A$. The last conclusion of \cref{lem:bt} is not stated in the reference, but follows immediately from the first conclusion and \cref{lem:tripling->AG}, below.

The purpose of the present note is to show that, for sufficiently large radii, a genuine small-doubling hypothesis is sufficient in \cref{lem:bt}. Indeed, we show this more generally for locally compact groups, in which context nothing weaker than small tripling was previously known to suffice.
\begin{theorem}\label{thm:doubling.lc}
Let $G$ be a locally compact group with Haar measure $\mu$, and let $K\ge1$. Suppose $S\subseteq G$ is a precompact symmetric open set containing the identity such that $\mu(S^{2n})\le K\mu(S^n)$ for some integer $n\ge8K^4$. Then $\mu(S^{3n})\le(2^{12}K^{24})^{10\cdot5^{4K^4}}\mu(S^n)$, and $S^{2n}$ is a precompact open $(2^{12}K^{24})^{30\cdot5^{4K^4}}$-approximate group.
\end{theorem}
Here and throughout, by a Haar measure we mean a \emph{left Haar} measure, so that $\mu(xA)=\mu(A)$ for all $A\subseteq G$ and $x\in G$.

A particular reason to be interested in locally compact groups is that the automorphism group of a vertex-transitive graph $\Gamma$ is a locally compact group with respect to the topology of pointwise convergence (see e.g. \cite[\S4]{ttTrof} for details). Given a vertex $o\in\Gamma$, the set $S=\{g\in\Aut(\Gamma):d(o,g(o))\le1\}$ is a compact open generating set for $\Aut(\Gamma)$; moreover, if we write $\beta_\Gamma(n)$ for the number of vertices in a ball of radius $n$ in $\Gamma$, and $\mu$ for the left Haar measure on $\Aut(\Gamma)$ normalised so that the stabiliser of $o$ has measure $1$, then we have $\beta_\Gamma(n)=\mu(S^n)$ for each $n\in\N$ \cite[Lemma 4.8]{ttTrof}. \cref{thm:doubling.lc} therefore gives rise to the following corollary.
\begin{corollary}\label{cor:trans}
Let $K\ge1$. Suppose $\Gamma$ is a locally finite vertex-transitive graph satisfying $\beta_\Gamma(2n)\le K\beta_\Gamma(n)$ for some integer $n\ge8K^4$. Then $\beta_\Gamma(3n)\le(2^{12}K^{24})^{10\cdot5^{4K^4}}\beta_\Gamma(n)$.
\end{corollary}
In the special case of discrete groups in \cref{thm:doubling.lc}, or equivalently Cayley graphs in \cref{cor:trans}, our argument gives slightly stronger bounds, as follows.
\begin{theorem}\label{thm:doubling=>tripling}
Let $K\ge1$. Suppose $S$ is a finite symmetric subset of a group containing the identity, and $|S^{2n}|\le K|S^n|$ for some integer $n\ge2K^2$. Then $|S^{3n}|\le K^3(3^9K^{18})^{10\cdot5^{K^2}}|S^n|$, and $S^{2n}$ is an $K^9(3^9K^{18})^{30\cdot5^{K^2}}$-approximate group.
\end{theorem}
We see no reason to believe that the bounds of \cref{thm:doubling.lc,thm:doubling=>tripling} are anything close to optimal.
